% Insert with % Insert with % Insert with % Insert with \input{paper/methods}; requires amsmath and amssymb.
\section{Methods}
\label{sec:methods}

\subsection{Data and analytical sample}
The analysis used households surveyed on one weekday in 2006 and the named fields in columns B--P of the supplied workbook. The unlabeled columns lacked definitions suitable for analysis. The response, $y_i$, was the recorded number of household trips (\texttt{nhhldtrip}). Data screening checked household identifiers, numerical ranges, categorical codes, positive expansion factors, and consistency among household counts. Licensed drivers and students were bounded by household size; the combined full-time and part-time worker count was also bounded by household size, and home workers by the total worker count.

The common eligibility rule excluded 33 households reporting ten or more vehicles and nine households with unknown dwelling type; the exclusions did not overlap. This left 93,265 of the 93,307 supplied households. The vehicle threshold restricts inference to households with fewer than ten reported vehicles; records coded as 99 vehicles were among those excluded. Zero-trip households remained eligible. Both modeling methods used the same sample.

The survey records trips for household members aged 11 and older, while household size includes younger children and the student count includes assumed students aged 6--10. The household counts therefore cover a wider age range than the trip outcome. The outcome measures trip frequency on the recorded weekday, regardless of purpose, mode, distance, or destination.

\subsection{Reproducible sample division and exploratory analysis}
Household identifiers were sorted and randomly permuted using NumPy's PCG64 generator with seed 42. The first rounded 30\% were assigned to testing and the remainder to training. The split preceded eligibility exclusions and had no stratification. Applying the common exclusions produced 65,287 training households and 27,978 test households. Sorting the identifiers made the partition reproducible across different source-row orders. The regression and cross-classification workflows regenerated the assignment and checked exact agreement with the EDA split saved in the training workbook.

Exploratory summaries used eligible training households. Distributions, quantiles, zero-trip shares, and group trip rates described the outcome and household characteristics. Pearson correlations assessed linear associations on the original count scale, particularly overlap among predictors entering OLS. Because the trip distribution has a long upper tail, its correlations were interpreted alongside group profiles and influence diagnostics \cite{deWinter2016}. The analysis also checked support for joint household categories and compared unweighted and expansion-weighted means. Display caps pooled sparse upper categories in figures and cross-classification; regression retained the original eligible counts.

The eligible sample covered 1,179 traffic zones, including 1,158 represented in training and 1,081 in testing. Because households from a zone could occur in both partitions, the test measures prediction within the 2006 survey setting. The fixed test partition remained outside coefficient estimation, information-criterion ranking, and estimation of cross-classification rates. Earlier workflow development involved exploration of this dataset; the resulting test assessment is therefore an internal validation exercise.

\subsection{Multiple linear regression}
\begin{samepage}
Multiple linear regression was used to describe the conditional association between household characteristics and the expected number of recorded trips. For household $i$, the working model was
\begin{equation}
 y_i = \beta_0 + \sum_{j=1}^{p-1} x_{ij}\beta_j + \varepsilon_i,
 \qquad \widehat y_i = \widehat\beta_0 + \sum_{j=1}^{p-1}x_{ij}\widehat\beta_j,
 \label{eq:mlr}
\end{equation}
\end{samepage}
where $x_{ij}$ is a predictor or constructed predictor term, $p$ includes the intercept, and $\varepsilon_i$ represents variation not captured by the specification. The model is linear in its coefficients when predictors are multiplied to represent interactions. Its predictions are conditional mean trip rates and need not be integers.

Let $P_i$, $V_i$, $L_i$, $F_i$, $Q_i$, $H_i$, and $S_i$ denote household persons, vehicles, licensed drivers, full-time workers, part-time workers, home workers, and students, respectively. Household size, vehicle count, and driver count were centered as $P_{ci}=P_i-3$, $V_{ci}=V_i-1$, and $L_{ci}=L_i-2$. In specifications with interactions, the coefficient of a centered count is a local slope at the reference values of the interacting counts. A one-unit change is evaluated from the full equation. These fixed reference values were chosen for this project. Centering leaves the fitted values of the corresponding complete specification unchanged.

The starting specification was
\begin{equation}
 y_i=\beta_0+\beta_P P_{ci}+\beta_V V_{ci}
       +\beta_F F_i+\beta_S S_i+\varepsilon_i.
 \label{eq:simple-model}
\end{equation}
Persons, workers, and students represent potential participation in daily activities, while vehicles represent access to motorized travel. The coefficients describe conditional associations within the observed households.

The candidate set extended the starting model to licensed drivers, alternative employment measures, dwelling type, region, and two interactions. Driver count indicates capacity for independent motorized travel; home working may change commuting and other travel. Dwelling type and region summarize residential context. Interactions let the driver and vehicle associations vary with household size.

The searched effects were selected from the following encompassing specification:
\begin{align}
 y_i={}&\beta_0+\beta_P P_{ci}+\beta_V V_{ci}+\beta_L L_{ci}
       +\beta_F F_i+\beta_Q Q_i+\beta_H H_i+\beta_S S_i \notag\\
      &+\boldsymbol{\gamma}^{\mathsf T}\mathbf d_i
       +\boldsymbol{\eta}^{\mathsf T}\mathbf r_i \notag\\
      &+\beta_{PL}P_{ci}L_{ci}+\beta_{PV}P_{ci}V_{ci}
       +\varepsilon_i.
 \label{eq:alternative-model}
\end{align}
Here $\mathbf d_i$ and $\mathbf r_i$ contain indicator variables for dwelling type and region. Their reference categories were House and Toronto. The other regions were Durham, York, and Peel, assigned from residential zone ranges. The search treated total workers, $F_i+Q_i$, as a substitute for the separate employment counts, avoiding redundant effects. It used the region categories to summarize location, while household identifiers, exact zone identifiers, survey administration fields, and the expansion factor remained outside the predictor set.

\subsubsection{Estimation and specification search}
All coefficients were estimated by unweighted ordinary least squares on the training households:
\begin{equation}
 \widehat{\boldsymbol\beta}
 =\underset{\boldsymbol\beta}{\operatorname{arg\,min}}
   \sum_{i\in\mathcal T}(y_i-\mathbf x_i^{\mathsf T}\boldsymbol\beta)^2,
 \label{eq:ols}
\end{equation}
where $\mathcal T$ denotes the training set and $\mathbf x_i$ includes the intercept. This gives each sampled household equal weight in fitting. The same training observations were used for every candidate model.

The search evaluated 1,040 valid effect subsets, including the intercept-only specification; household size was optional. Categorical variables entered as whole factors, and interactions required their supporting main effects. Candidates were scored using working Gaussian OLS likelihoods:
\begin{equation}
 \mathrm{AIC}=-2\widehat\ell+2p,
 \qquad
 \mathrm{BIC}=-2\widehat\ell+p\log(n_{\mathrm{tr}}),
 \label{eq:ic}
\end{equation}
where $n_{\mathrm{tr}}=|\mathcal T|$ and $\widehat\ell$ is the maximized working log-likelihood. The penalty counts the regression coefficients, including the intercept, in $p$; a common variance-parameter penalty leaves the ranking unchanged. The working Gaussian likelihood provided a comparison rule, while the observed trip-count distribution informed the residual and prediction checks.

Models were examined in ascending BIC order. The selected alternative was the first model for which every removable whole effect had a zone-clustered Wald test $p$-value below 0.01. Only main effects required by retained interactions were protected from removal. The AIC leaders and the best additive specification were retained for comparison. The squared terms were removed after the first evaluation because their parabolic form made predictions fall at high vehicle counts and produced negative predictions. Survey weekday was subsequently removed from the candidate effects to simplify the specification. The revised search used training data; its test results constitute internal validation after this specification revision.

Coefficient uncertainty was estimated using a sandwich covariance clustered by residential traffic zone, with the finite-sample correction
$G/(G-1)\,(n_{\mathrm{tr}}-1)/(n_{\mathrm{tr}}-p)$, where $G=1,158$ training zones. Coefficient intervals used a $t$ reference distribution with $G-1$ degrees of freedom; whole-factor tests used the corresponding clustered Wald $F$ statistic. This permits correlated errors within zones and unequal error variances, while retaining the OLS coefficient estimates. The reported 99\% intervals and tests for the selected alternative are exploratory because they do not adjust for searching among specifications.

Training goodness of fit was summarized by
\begin{equation}
 R^2=1-\frac{\sum_{i\in\mathcal T}(y_i-\widehat y_i)^2}
                 {\sum_{i\in\mathcal T}(y_i-\overline y_{\mathcal T})^2},
 \qquad
 \overline R^2=1-(1-R^2)\frac{n_{\mathrm{tr}}-1}{n_{\mathrm{tr}}-p}.
\end{equation}
Residual plots, a normal quantile plot, residual skewness and kurtosis, variance-inflation factors, leverage, Cook's distance, and externally studentized residuals were examined. Leverage above $3p/n_{\mathrm{tr}}$, Cook's distance above $4/n_{\mathrm{tr}}$, and absolute studentized residuals above three flagged observations for review. The robust Breusch--Pagan diagnostic provided an independent-observation screen for heteroscedasticity. Predicted trip curves at fixed household profiles were examined across the observed vehicle range together with exact-count training support. Negative predictions remained in the evaluation so that errors represented the fitted OLS model.

\subsection{Cross-classification}
The cross-classification model assigned an average trip rate to households sharing a size and licensed-driver category, an established trip-generation approach \cite{Moeckel2017}. The primary size groups were 1, 2, 3, 4, and 5 or more persons; driver groups were 0, 1, 2, and 3 or more. Household size represents potential activity demand, while driver availability describes one aspect of mobility access. Their joint categories allow mean trip rates to vary across combinations. Pooling the upper categories increases sample support in the tails.

For size group $a$ and driver group $b$, let $\mathcal T_{ab}$ be the corresponding training households and $n_{ab}=|\mathcal T_{ab}|$. The cell rate was the unweighted mean
\begin{equation}
 \widehat r_{ab}=\frac{1}{n_{ab}}\sum_{i\in\mathcal T_{ab}}y_i.
\end{equation}
A feasible cell with fewer than 30 training households used the training mean for its size group, $\widehat r_a$. Thus, the prediction for a feasible test household was
\begin{equation}
 \widehat y_i^{\mathrm{CC}}=
 \begin{cases}
 \widehat r_{ab}, & n_{ab}\geq30,\\
 \widehat r_a, & n_{ab}<30.
 \end{cases}
\end{equation}
The threshold of 30 set a minimum cell count for direct rate estimation. Precision still varied with cell size and within-cell variation. Structurally impossible combinations, such as two licensed drivers in a one-person household, had blank rates and heatmap cells. Sparse-cell counts and fallback predictions applied to feasible combinations. Tables and the heatmap reported the household count supporting each observed rate.

Two alternative boundary schemes were evaluated: size capped at 4 or more with drivers capped at 3 or more, and size capped at 5 or more with drivers capped at 4 or more. All rates and fallback values were fitted on training households. Test results for these alternatives served as boundary sensitivity checks. The primary boundaries were fixed from the training-data support assessment.

\subsection{Household prediction assessment}
Regression and cross-classification were evaluated on the same $n_{\mathrm{te}}=27,978$ eligible test households. A benchmark assigned every test household the unweighted training mean, $\widehat y_i^{\mathrm{mean}}=\overline y_{\mathcal T}$. A regression using only household size and licensed drivers gave a common-predictor comparison with cross-classification; it retained exact counts where cross-classification grouped upper counts. The simple four-predictor regression, best additive alternative, and selected regression showed how performance changed across broader specifications. Both principal methods were also assessed within the same size and driver groups.

For prediction error $e_i=\widehat y_i-y_i$, the household measures were
\begin{align}
 \mathrm{MAE}&=\frac{1}{n_{\mathrm{te}}}\sum_{i\in\mathcal E}|e_i|,\\
 \mathrm{RMSE}&=\sqrt{\frac{1}{n_{\mathrm{te}}}\sum_{i\in\mathcal E}e_i^2},\\
 \mathrm{ME}&=\frac{1}{n_{\mathrm{te}}}\sum_{i\in\mathcal E}e_i,
\end{align}
where $\mathcal E$ is the test set. MAE gives the average absolute error in trips per household. RMSE gives greater weight to large errors. ME is the signed mean error in trips: a positive value indicates overprediction. A small ME can coexist with substantial household errors because positive and negative errors cancel. Observed-versus-predicted plots and results for selected household groups supplemented these overall measures.

The expansion factor $w_i=\texttt{expf}_i$ was used for an additional weighted assessment, without re-estimating either model. Weighted MAE and ME replaced $1/n_{\mathrm{te}}$ by normalized weights $w_i/\sum_{j\in\mathcal E}w_j$; weighted RMSE was
\begin{equation}
 \mathrm{RMSE}_{w}=\sqrt{\frac{\sum_{i\in\mathcal E}w_i e_i^2}
                                  {\sum_{i\in\mathcal E}w_i}}.
\end{equation}

\subsection{Traffic-zone aggregation and scope of inference}
For each zone $z$ represented in the test set, observed and predicted sample totals were summed over the same records $\mathcal E_z$:
\begin{equation}
 Y_z=\sum_{i\in\mathcal E_z}y_i,
 \qquad \widehat Y_z=\sum_{i\in\mathcal E_z}\widehat y_i.
\end{equation}
Expansion-weighted sums were calculated consistently as
\begin{equation}
 Y_z^{w}=\sum_{i\in\mathcal E_z}w_i y_i,
 \qquad \widehat Y_z^{w}=\sum_{i\in\mathcal E_z}w_i\widehat y_i.
\end{equation}
Zone MAE, RMSE, and ME used the corresponding differences, with each of the 1,081 evaluated zones receiving equal weight in these summaries. Zones without test households were not included. These measures and zone agreement plots assess how household errors accumulate or cancel when aggregated.

The weighted sums represent the contributions of the test records under their original survey expansion factors. They are sample-based zone aggregates because the analysis used neither an inverse test-selection adjustment nor population recalibration. Both model comparisons use the same records and weighting convention. The results describe associations and predictive performance within this historical survey sample. Forecasts for another period or geography would require separate validation.
; requires amsmath and amssymb.
\section{Methods}
\label{sec:methods}

\subsection{Data and analytical sample}
The analysis used households surveyed on one weekday in 2006 and the named fields in columns B--P of the supplied workbook. The unlabeled columns lacked definitions suitable for analysis. The response, $y_i$, was the recorded number of household trips (\texttt{nhhldtrip}). Data screening checked household identifiers, numerical ranges, categorical codes, positive expansion factors, and consistency among household counts. Licensed drivers and students were bounded by household size; the combined full-time and part-time worker count was also bounded by household size, and home workers by the total worker count.

The common eligibility rule excluded 33 households reporting ten or more vehicles and nine households with unknown dwelling type; the exclusions did not overlap. This left 93,265 of the 93,307 supplied households. The vehicle threshold restricts inference to households with fewer than ten reported vehicles; records coded as 99 vehicles were among those excluded. Zero-trip households remained eligible. Both modeling methods used the same sample.

The survey records trips for household members aged 11 and older, while household size includes younger children and the student count includes assumed students aged 6--10. The household counts therefore cover a wider age range than the trip outcome. The outcome measures trip frequency on the recorded weekday, regardless of purpose, mode, distance, or destination.

\subsection{Reproducible sample division and exploratory analysis}
Household identifiers were sorted and randomly permuted using NumPy's PCG64 generator with seed 42. The first rounded 30\% were assigned to testing and the remainder to training. The split preceded eligibility exclusions and had no stratification. Applying the common exclusions produced 65,287 training households and 27,978 test households. Sorting the identifiers made the partition reproducible across different source-row orders. The regression and cross-classification workflows regenerated the assignment and checked exact agreement with the EDA split saved in the training workbook.

Exploratory summaries used eligible training households. Distributions, quantiles, zero-trip shares, and group trip rates described the outcome and household characteristics. Pearson correlations assessed linear associations on the original count scale, particularly overlap among predictors entering OLS. Because the trip distribution has a long upper tail, its correlations were interpreted alongside group profiles and influence diagnostics \cite{deWinter2016}. The analysis also checked support for joint household categories and compared unweighted and expansion-weighted means. Display caps pooled sparse upper categories in figures and cross-classification; regression retained the original eligible counts.

The eligible sample covered 1,179 traffic zones, including 1,158 represented in training and 1,081 in testing. Because households from a zone could occur in both partitions, the test measures prediction within the 2006 survey setting. The fixed test partition remained outside coefficient estimation, information-criterion ranking, and estimation of cross-classification rates. Earlier workflow development involved exploration of this dataset; the resulting test assessment is therefore an internal validation exercise.

\subsection{Multiple linear regression}
\begin{samepage}
Multiple linear regression was used to describe the conditional association between household characteristics and the expected number of recorded trips. For household $i$, the working model was
\begin{equation}
 y_i = \beta_0 + \sum_{j=1}^{p-1} x_{ij}\beta_j + \varepsilon_i,
 \qquad \widehat y_i = \widehat\beta_0 + \sum_{j=1}^{p-1}x_{ij}\widehat\beta_j,
 \label{eq:mlr}
\end{equation}
\end{samepage}
where $x_{ij}$ is a predictor or constructed predictor term, $p$ includes the intercept, and $\varepsilon_i$ represents variation not captured by the specification. The model is linear in its coefficients when predictors are multiplied to represent interactions. Its predictions are conditional mean trip rates and need not be integers.

Let $P_i$, $V_i$, $L_i$, $F_i$, $Q_i$, $H_i$, and $S_i$ denote household persons, vehicles, licensed drivers, full-time workers, part-time workers, home workers, and students, respectively. Household size, vehicle count, and driver count were centered as $P_{ci}=P_i-3$, $V_{ci}=V_i-1$, and $L_{ci}=L_i-2$. In specifications with interactions, the coefficient of a centered count is a local slope at the reference values of the interacting counts. A one-unit change is evaluated from the full equation. These fixed reference values were chosen for this project. Centering leaves the fitted values of the corresponding complete specification unchanged.

The starting specification was
\begin{equation}
 y_i=\beta_0+\beta_P P_{ci}+\beta_V V_{ci}
       +\beta_F F_i+\beta_S S_i+\varepsilon_i.
 \label{eq:simple-model}
\end{equation}
Persons, workers, and students represent potential participation in daily activities, while vehicles represent access to motorized travel. The coefficients describe conditional associations within the observed households.

The candidate set extended the starting model to licensed drivers, alternative employment measures, dwelling type, region, and two interactions. Driver count indicates capacity for independent motorized travel; home working may change commuting and other travel. Dwelling type and region summarize residential context. Interactions let the driver and vehicle associations vary with household size.

The searched effects were selected from the following encompassing specification:
\begin{align}
 y_i={}&\beta_0+\beta_P P_{ci}+\beta_V V_{ci}+\beta_L L_{ci}
       +\beta_F F_i+\beta_Q Q_i+\beta_H H_i+\beta_S S_i \notag\\
      &+\boldsymbol{\gamma}^{\mathsf T}\mathbf d_i
       +\boldsymbol{\eta}^{\mathsf T}\mathbf r_i \notag\\
      &+\beta_{PL}P_{ci}L_{ci}+\beta_{PV}P_{ci}V_{ci}
       +\varepsilon_i.
 \label{eq:alternative-model}
\end{align}
Here $\mathbf d_i$ and $\mathbf r_i$ contain indicator variables for dwelling type and region. Their reference categories were House and Toronto. The other regions were Durham, York, and Peel, assigned from residential zone ranges. The search treated total workers, $F_i+Q_i$, as a substitute for the separate employment counts, avoiding redundant effects. It used the region categories to summarize location, while household identifiers, exact zone identifiers, survey administration fields, and the expansion factor remained outside the predictor set.

\subsubsection{Estimation and specification search}
All coefficients were estimated by unweighted ordinary least squares on the training households:
\begin{equation}
 \widehat{\boldsymbol\beta}
 =\underset{\boldsymbol\beta}{\operatorname{arg\,min}}
   \sum_{i\in\mathcal T}(y_i-\mathbf x_i^{\mathsf T}\boldsymbol\beta)^2,
 \label{eq:ols}
\end{equation}
where $\mathcal T$ denotes the training set and $\mathbf x_i$ includes the intercept. This gives each sampled household equal weight in fitting. The same training observations were used for every candidate model.

The search evaluated 1,040 valid effect subsets, including the intercept-only specification; household size was optional. Categorical variables entered as whole factors, and interactions required their supporting main effects. Candidates were scored using working Gaussian OLS likelihoods:
\begin{equation}
 \mathrm{AIC}=-2\widehat\ell+2p,
 \qquad
 \mathrm{BIC}=-2\widehat\ell+p\log(n_{\mathrm{tr}}),
 \label{eq:ic}
\end{equation}
where $n_{\mathrm{tr}}=|\mathcal T|$ and $\widehat\ell$ is the maximized working log-likelihood. The penalty counts the regression coefficients, including the intercept, in $p$; a common variance-parameter penalty leaves the ranking unchanged. The working Gaussian likelihood provided a comparison rule, while the observed trip-count distribution informed the residual and prediction checks.

Models were examined in ascending BIC order. The selected alternative was the first model for which every removable whole effect had a zone-clustered Wald test $p$-value below 0.01. Only main effects required by retained interactions were protected from removal. The AIC leaders and the best additive specification were retained for comparison. The squared terms were removed after the first evaluation because their parabolic form made predictions fall at high vehicle counts and produced negative predictions. Survey weekday was subsequently removed from the candidate effects to simplify the specification. The revised search used training data; its test results constitute internal validation after this specification revision.

Coefficient uncertainty was estimated using a sandwich covariance clustered by residential traffic zone, with the finite-sample correction
$G/(G-1)\,(n_{\mathrm{tr}}-1)/(n_{\mathrm{tr}}-p)$, where $G=1,158$ training zones. Coefficient intervals used a $t$ reference distribution with $G-1$ degrees of freedom; whole-factor tests used the corresponding clustered Wald $F$ statistic. This permits correlated errors within zones and unequal error variances, while retaining the OLS coefficient estimates. The reported 99\% intervals and tests for the selected alternative are exploratory because they do not adjust for searching among specifications.

Training goodness of fit was summarized by
\begin{equation}
 R^2=1-\frac{\sum_{i\in\mathcal T}(y_i-\widehat y_i)^2}
                 {\sum_{i\in\mathcal T}(y_i-\overline y_{\mathcal T})^2},
 \qquad
 \overline R^2=1-(1-R^2)\frac{n_{\mathrm{tr}}-1}{n_{\mathrm{tr}}-p}.
\end{equation}
Residual plots, a normal quantile plot, residual skewness and kurtosis, variance-inflation factors, leverage, Cook's distance, and externally studentized residuals were examined. Leverage above $3p/n_{\mathrm{tr}}$, Cook's distance above $4/n_{\mathrm{tr}}$, and absolute studentized residuals above three flagged observations for review. The robust Breusch--Pagan diagnostic provided an independent-observation screen for heteroscedasticity. Predicted trip curves at fixed household profiles were examined across the observed vehicle range together with exact-count training support. Negative predictions remained in the evaluation so that errors represented the fitted OLS model.

\subsection{Cross-classification}
The cross-classification model assigned an average trip rate to households sharing a size and licensed-driver category, an established trip-generation approach \cite{Moeckel2017}. The primary size groups were 1, 2, 3, 4, and 5 or more persons; driver groups were 0, 1, 2, and 3 or more. Household size represents potential activity demand, while driver availability describes one aspect of mobility access. Their joint categories allow mean trip rates to vary across combinations. Pooling the upper categories increases sample support in the tails.

For size group $a$ and driver group $b$, let $\mathcal T_{ab}$ be the corresponding training households and $n_{ab}=|\mathcal T_{ab}|$. The cell rate was the unweighted mean
\begin{equation}
 \widehat r_{ab}=\frac{1}{n_{ab}}\sum_{i\in\mathcal T_{ab}}y_i.
\end{equation}
A feasible cell with fewer than 30 training households used the training mean for its size group, $\widehat r_a$. Thus, the prediction for a feasible test household was
\begin{equation}
 \widehat y_i^{\mathrm{CC}}=
 \begin{cases}
 \widehat r_{ab}, & n_{ab}\geq30,\\
 \widehat r_a, & n_{ab}<30.
 \end{cases}
\end{equation}
The threshold of 30 set a minimum cell count for direct rate estimation. Precision still varied with cell size and within-cell variation. Structurally impossible combinations, such as two licensed drivers in a one-person household, had blank rates and heatmap cells. Sparse-cell counts and fallback predictions applied to feasible combinations. Tables and the heatmap reported the household count supporting each observed rate.

Two alternative boundary schemes were evaluated: size capped at 4 or more with drivers capped at 3 or more, and size capped at 5 or more with drivers capped at 4 or more. All rates and fallback values were fitted on training households. Test results for these alternatives served as boundary sensitivity checks. The primary boundaries were fixed from the training-data support assessment.

\subsection{Household prediction assessment}
Regression and cross-classification were evaluated on the same $n_{\mathrm{te}}=27,978$ eligible test households. A benchmark assigned every test household the unweighted training mean, $\widehat y_i^{\mathrm{mean}}=\overline y_{\mathcal T}$. A regression using only household size and licensed drivers gave a common-predictor comparison with cross-classification; it retained exact counts where cross-classification grouped upper counts. The simple four-predictor regression, best additive alternative, and selected regression showed how performance changed across broader specifications. Both principal methods were also assessed within the same size and driver groups.

For prediction error $e_i=\widehat y_i-y_i$, the household measures were
\begin{align}
 \mathrm{MAE}&=\frac{1}{n_{\mathrm{te}}}\sum_{i\in\mathcal E}|e_i|,\\
 \mathrm{RMSE}&=\sqrt{\frac{1}{n_{\mathrm{te}}}\sum_{i\in\mathcal E}e_i^2},\\
 \mathrm{ME}&=\frac{1}{n_{\mathrm{te}}}\sum_{i\in\mathcal E}e_i,
\end{align}
where $\mathcal E$ is the test set. MAE gives the average absolute error in trips per household. RMSE gives greater weight to large errors. ME is the signed mean error in trips: a positive value indicates overprediction. A small ME can coexist with substantial household errors because positive and negative errors cancel. Observed-versus-predicted plots and results for selected household groups supplemented these overall measures.

The expansion factor $w_i=\texttt{expf}_i$ was used for an additional weighted assessment, without re-estimating either model. Weighted MAE and ME replaced $1/n_{\mathrm{te}}$ by normalized weights $w_i/\sum_{j\in\mathcal E}w_j$; weighted RMSE was
\begin{equation}
 \mathrm{RMSE}_{w}=\sqrt{\frac{\sum_{i\in\mathcal E}w_i e_i^2}
                                  {\sum_{i\in\mathcal E}w_i}}.
\end{equation}

\subsection{Traffic-zone aggregation and scope of inference}
For each zone $z$ represented in the test set, observed and predicted sample totals were summed over the same records $\mathcal E_z$:
\begin{equation}
 Y_z=\sum_{i\in\mathcal E_z}y_i,
 \qquad \widehat Y_z=\sum_{i\in\mathcal E_z}\widehat y_i.
\end{equation}
Expansion-weighted sums were calculated consistently as
\begin{equation}
 Y_z^{w}=\sum_{i\in\mathcal E_z}w_i y_i,
 \qquad \widehat Y_z^{w}=\sum_{i\in\mathcal E_z}w_i\widehat y_i.
\end{equation}
Zone MAE, RMSE, and ME used the corresponding differences, with each of the 1,081 evaluated zones receiving equal weight in these summaries. Zones without test households were not included. These measures and zone agreement plots assess how household errors accumulate or cancel when aggregated.

The weighted sums represent the contributions of the test records under their original survey expansion factors. They are sample-based zone aggregates because the analysis used neither an inverse test-selection adjustment nor population recalibration. Both model comparisons use the same records and weighting convention. The results describe associations and predictive performance within this historical survey sample. Forecasts for another period or geography would require separate validation.
; requires amsmath and amssymb.
\section{Methods}
\label{sec:methods}

\subsection{Data and analytical sample}
The analysis used households surveyed on one weekday in 2006 and the named fields in columns B--P of the supplied workbook. The unlabeled columns lacked definitions suitable for analysis. The response, $y_i$, was the recorded number of household trips (\texttt{nhhldtrip}). Data screening checked household identifiers, numerical ranges, categorical codes, positive expansion factors, and consistency among household counts. Licensed drivers and students were bounded by household size; the combined full-time and part-time worker count was also bounded by household size, and home workers by the total worker count.

The common eligibility rule excluded 33 households reporting ten or more vehicles and nine households with unknown dwelling type; the exclusions did not overlap. This left 93,265 of the 93,307 supplied households. The vehicle threshold restricts inference to households with fewer than ten reported vehicles; records coded as 99 vehicles were among those excluded. Zero-trip households remained eligible. Both modeling methods used the same sample.

The survey records trips for household members aged 11 and older, while household size includes younger children and the student count includes assumed students aged 6--10. The household counts therefore cover a wider age range than the trip outcome. The outcome measures trip frequency on the recorded weekday, regardless of purpose, mode, distance, or destination.

\subsection{Reproducible sample division and exploratory analysis}
Household identifiers were sorted and randomly permuted using NumPy's PCG64 generator with seed 42. The first rounded 30\% were assigned to testing and the remainder to training. The split preceded eligibility exclusions and had no stratification. Applying the common exclusions produced 65,287 training households and 27,978 test households. Sorting the identifiers made the partition reproducible across different source-row orders. The regression and cross-classification workflows regenerated the assignment and checked exact agreement with the EDA split saved in the training workbook.

Exploratory summaries used eligible training households. Distributions, quantiles, zero-trip shares, and group trip rates described the outcome and household characteristics. Pearson correlations assessed linear associations on the original count scale, particularly overlap among predictors entering OLS. Because the trip distribution has a long upper tail, its correlations were interpreted alongside group profiles and influence diagnostics \cite{deWinter2016}. The analysis also checked support for joint household categories and compared unweighted and expansion-weighted means. Display caps pooled sparse upper categories in figures and cross-classification; regression retained the original eligible counts.

The eligible sample covered 1,179 traffic zones, including 1,158 represented in training and 1,081 in testing. Because households from a zone could occur in both partitions, the test measures prediction within the 2006 survey setting. The fixed test partition remained outside coefficient estimation, information-criterion ranking, and estimation of cross-classification rates. Earlier workflow development involved exploration of this dataset; the resulting test assessment is therefore an internal validation exercise.

\subsection{Multiple linear regression}
\begin{samepage}
Multiple linear regression was used to describe the conditional association between household characteristics and the expected number of recorded trips. For household $i$, the working model was
\begin{equation}
 y_i = \beta_0 + \sum_{j=1}^{p-1} x_{ij}\beta_j + \varepsilon_i,
 \qquad \widehat y_i = \widehat\beta_0 + \sum_{j=1}^{p-1}x_{ij}\widehat\beta_j,
 \label{eq:mlr}
\end{equation}
\end{samepage}
where $x_{ij}$ is a predictor or constructed predictor term, $p$ includes the intercept, and $\varepsilon_i$ represents variation not captured by the specification. The model is linear in its coefficients when predictors are multiplied to represent interactions. Its predictions are conditional mean trip rates and need not be integers.

Let $P_i$, $V_i$, $L_i$, $F_i$, $Q_i$, $H_i$, and $S_i$ denote household persons, vehicles, licensed drivers, full-time workers, part-time workers, home workers, and students, respectively. Household size, vehicle count, and driver count were centered as $P_{ci}=P_i-3$, $V_{ci}=V_i-1$, and $L_{ci}=L_i-2$. In specifications with interactions, the coefficient of a centered count is a local slope at the reference values of the interacting counts. A one-unit change is evaluated from the full equation. These fixed reference values were chosen for this project. Centering leaves the fitted values of the corresponding complete specification unchanged.

The starting specification was
\begin{equation}
 y_i=\beta_0+\beta_P P_{ci}+\beta_V V_{ci}
       +\beta_F F_i+\beta_S S_i+\varepsilon_i.
 \label{eq:simple-model}
\end{equation}
Persons, workers, and students represent potential participation in daily activities, while vehicles represent access to motorized travel. The coefficients describe conditional associations within the observed households.

The candidate set extended the starting model to licensed drivers, alternative employment measures, dwelling type, region, and two interactions. Driver count indicates capacity for independent motorized travel; home working may change commuting and other travel. Dwelling type and region summarize residential context. Interactions let the driver and vehicle associations vary with household size.

The searched effects were selected from the following encompassing specification:
\begin{align}
 y_i={}&\beta_0+\beta_P P_{ci}+\beta_V V_{ci}+\beta_L L_{ci}
       +\beta_F F_i+\beta_Q Q_i+\beta_H H_i+\beta_S S_i \notag\\
      &+\boldsymbol{\gamma}^{\mathsf T}\mathbf d_i
       +\boldsymbol{\eta}^{\mathsf T}\mathbf r_i \notag\\
      &+\beta_{PL}P_{ci}L_{ci}+\beta_{PV}P_{ci}V_{ci}
       +\varepsilon_i.
 \label{eq:alternative-model}
\end{align}
Here $\mathbf d_i$ and $\mathbf r_i$ contain indicator variables for dwelling type and region. Their reference categories were House and Toronto. The other regions were Durham, York, and Peel, assigned from residential zone ranges. The search treated total workers, $F_i+Q_i$, as a substitute for the separate employment counts, avoiding redundant effects. It used the region categories to summarize location, while household identifiers, exact zone identifiers, survey administration fields, and the expansion factor remained outside the predictor set.

\subsubsection{Estimation and specification search}
All coefficients were estimated by unweighted ordinary least squares on the training households:
\begin{equation}
 \widehat{\boldsymbol\beta}
 =\underset{\boldsymbol\beta}{\operatorname{arg\,min}}
   \sum_{i\in\mathcal T}(y_i-\mathbf x_i^{\mathsf T}\boldsymbol\beta)^2,
 \label{eq:ols}
\end{equation}
where $\mathcal T$ denotes the training set and $\mathbf x_i$ includes the intercept. This gives each sampled household equal weight in fitting. The same training observations were used for every candidate model.

The search evaluated 1,040 valid effect subsets, including the intercept-only specification; household size was optional. Categorical variables entered as whole factors, and interactions required their supporting main effects. Candidates were scored using working Gaussian OLS likelihoods:
\begin{equation}
 \mathrm{AIC}=-2\widehat\ell+2p,
 \qquad
 \mathrm{BIC}=-2\widehat\ell+p\log(n_{\mathrm{tr}}),
 \label{eq:ic}
\end{equation}
where $n_{\mathrm{tr}}=|\mathcal T|$ and $\widehat\ell$ is the maximized working log-likelihood. The penalty counts the regression coefficients, including the intercept, in $p$; a common variance-parameter penalty leaves the ranking unchanged. The working Gaussian likelihood provided a comparison rule, while the observed trip-count distribution informed the residual and prediction checks.

Models were examined in ascending BIC order. The selected alternative was the first model for which every removable whole effect had a zone-clustered Wald test $p$-value below 0.01. Only main effects required by retained interactions were protected from removal. The AIC leaders and the best additive specification were retained for comparison. The squared terms were removed after the first evaluation because their parabolic form made predictions fall at high vehicle counts and produced negative predictions. Survey weekday was subsequently removed from the candidate effects to simplify the specification. The revised search used training data; its test results constitute internal validation after this specification revision.

Coefficient uncertainty was estimated using a sandwich covariance clustered by residential traffic zone, with the finite-sample correction
$G/(G-1)\,(n_{\mathrm{tr}}-1)/(n_{\mathrm{tr}}-p)$, where $G=1,158$ training zones. Coefficient intervals used a $t$ reference distribution with $G-1$ degrees of freedom; whole-factor tests used the corresponding clustered Wald $F$ statistic. This permits correlated errors within zones and unequal error variances, while retaining the OLS coefficient estimates. The reported 99\% intervals and tests for the selected alternative are exploratory because they do not adjust for searching among specifications.

Training goodness of fit was summarized by
\begin{equation}
 R^2=1-\frac{\sum_{i\in\mathcal T}(y_i-\widehat y_i)^2}
                 {\sum_{i\in\mathcal T}(y_i-\overline y_{\mathcal T})^2},
 \qquad
 \overline R^2=1-(1-R^2)\frac{n_{\mathrm{tr}}-1}{n_{\mathrm{tr}}-p}.
\end{equation}
Residual plots, a normal quantile plot, residual skewness and kurtosis, variance-inflation factors, leverage, Cook's distance, and externally studentized residuals were examined. Leverage above $3p/n_{\mathrm{tr}}$, Cook's distance above $4/n_{\mathrm{tr}}$, and absolute studentized residuals above three flagged observations for review. The robust Breusch--Pagan diagnostic provided an independent-observation screen for heteroscedasticity. Predicted trip curves at fixed household profiles were examined across the observed vehicle range together with exact-count training support. Negative predictions remained in the evaluation so that errors represented the fitted OLS model.

\subsection{Cross-classification}
The cross-classification model assigned an average trip rate to households sharing a size and licensed-driver category, an established trip-generation approach \cite{Moeckel2017}. The primary size groups were 1, 2, 3, 4, and 5 or more persons; driver groups were 0, 1, 2, and 3 or more. Household size represents potential activity demand, while driver availability describes one aspect of mobility access. Their joint categories allow mean trip rates to vary across combinations. Pooling the upper categories increases sample support in the tails.

For size group $a$ and driver group $b$, let $\mathcal T_{ab}$ be the corresponding training households and $n_{ab}=|\mathcal T_{ab}|$. The cell rate was the unweighted mean
\begin{equation}
 \widehat r_{ab}=\frac{1}{n_{ab}}\sum_{i\in\mathcal T_{ab}}y_i.
\end{equation}
A feasible cell with fewer than 30 training households used the training mean for its size group, $\widehat r_a$. Thus, the prediction for a feasible test household was
\begin{equation}
 \widehat y_i^{\mathrm{CC}}=
 \begin{cases}
 \widehat r_{ab}, & n_{ab}\geq30,\\
 \widehat r_a, & n_{ab}<30.
 \end{cases}
\end{equation}
The threshold of 30 set a minimum cell count for direct rate estimation. Precision still varied with cell size and within-cell variation. Structurally impossible combinations, such as two licensed drivers in a one-person household, had blank rates and heatmap cells. Sparse-cell counts and fallback predictions applied to feasible combinations. Tables and the heatmap reported the household count supporting each observed rate.

Two alternative boundary schemes were evaluated: size capped at 4 or more with drivers capped at 3 or more, and size capped at 5 or more with drivers capped at 4 or more. All rates and fallback values were fitted on training households. Test results for these alternatives served as boundary sensitivity checks. The primary boundaries were fixed from the training-data support assessment.

\subsection{Household prediction assessment}
Regression and cross-classification were evaluated on the same $n_{\mathrm{te}}=27,978$ eligible test households. A benchmark assigned every test household the unweighted training mean, $\widehat y_i^{\mathrm{mean}}=\overline y_{\mathcal T}$. A regression using only household size and licensed drivers gave a common-predictor comparison with cross-classification; it retained exact counts where cross-classification grouped upper counts. The simple four-predictor regression, best additive alternative, and selected regression showed how performance changed across broader specifications. Both principal methods were also assessed within the same size and driver groups.

For prediction error $e_i=\widehat y_i-y_i$, the household measures were
\begin{align}
 \mathrm{MAE}&=\frac{1}{n_{\mathrm{te}}}\sum_{i\in\mathcal E}|e_i|,\\
 \mathrm{RMSE}&=\sqrt{\frac{1}{n_{\mathrm{te}}}\sum_{i\in\mathcal E}e_i^2},\\
 \mathrm{ME}&=\frac{1}{n_{\mathrm{te}}}\sum_{i\in\mathcal E}e_i,
\end{align}
where $\mathcal E$ is the test set. MAE gives the average absolute error in trips per household. RMSE gives greater weight to large errors. ME is the signed mean error in trips: a positive value indicates overprediction. A small ME can coexist with substantial household errors because positive and negative errors cancel. Observed-versus-predicted plots and results for selected household groups supplemented these overall measures.

The expansion factor $w_i=\texttt{expf}_i$ was used for an additional weighted assessment, without re-estimating either model. Weighted MAE and ME replaced $1/n_{\mathrm{te}}$ by normalized weights $w_i/\sum_{j\in\mathcal E}w_j$; weighted RMSE was
\begin{equation}
 \mathrm{RMSE}_{w}=\sqrt{\frac{\sum_{i\in\mathcal E}w_i e_i^2}
                                  {\sum_{i\in\mathcal E}w_i}}.
\end{equation}

\subsection{Traffic-zone aggregation and scope of inference}
For each zone $z$ represented in the test set, observed and predicted sample totals were summed over the same records $\mathcal E_z$:
\begin{equation}
 Y_z=\sum_{i\in\mathcal E_z}y_i,
 \qquad \widehat Y_z=\sum_{i\in\mathcal E_z}\widehat y_i.
\end{equation}
Expansion-weighted sums were calculated consistently as
\begin{equation}
 Y_z^{w}=\sum_{i\in\mathcal E_z}w_i y_i,
 \qquad \widehat Y_z^{w}=\sum_{i\in\mathcal E_z}w_i\widehat y_i.
\end{equation}
Zone MAE, RMSE, and ME used the corresponding differences, with each of the 1,081 evaluated zones receiving equal weight in these summaries. Zones without test households were not included. These measures and zone agreement plots assess how household errors accumulate or cancel when aggregated.

The weighted sums represent the contributions of the test records under their original survey expansion factors. They are sample-based zone aggregates because the analysis used neither an inverse test-selection adjustment nor population recalibration. Both model comparisons use the same records and weighting convention. The results describe associations and predictive performance within this historical survey sample. Forecasts for another period or geography would require separate validation.
; requires amsmath and amssymb.
\section{Methods}
\label{sec:methods}

\subsection{Data and analytical sample}
The analysis used households surveyed on one weekday in 2006 and the named fields in columns B--P of the supplied workbook. The unlabeled columns lacked definitions suitable for analysis. The response, $y_i$, was the recorded number of household trips (\texttt{nhhldtrip}). Data screening checked household identifiers, numerical ranges, categorical codes, positive expansion factors, and consistency among household counts. Licensed drivers and students were bounded by household size; the combined full-time and part-time worker count was also bounded by household size, and home workers by the total worker count.

The common eligibility rule excluded 33 households reporting ten or more vehicles and nine households with unknown dwelling type; the exclusions did not overlap. This left 93,265 of the 93,307 supplied households. The vehicle threshold restricts inference to households with fewer than ten reported vehicles; records coded as 99 vehicles were among those excluded. Zero-trip households remained eligible. Both modeling methods used the same sample.

The survey records trips for household members aged 11 and older, while household size includes younger children and the student count includes assumed students aged 6--10. The household counts therefore cover a wider age range than the trip outcome. The outcome measures trip frequency on the recorded weekday, regardless of purpose, mode, distance, or destination.

\subsection{Reproducible sample division and exploratory analysis}
Household identifiers were sorted and randomly permuted using NumPy's PCG64 generator with seed 42. The first rounded 30\% were assigned to testing and the remainder to training. The split preceded eligibility exclusions and had no stratification. Applying the common exclusions produced 65,287 training households and 27,978 test households. Sorting the identifiers made the partition reproducible across different source-row orders. The regression and cross-classification workflows regenerated the assignment and checked exact agreement with the EDA split saved in the training workbook.

Exploratory summaries used eligible training households. Distributions, quantiles, zero-trip shares, and group trip rates described the outcome and household characteristics. Pearson correlations assessed linear associations on the original count scale, particularly overlap among predictors entering OLS. Because the trip distribution has a long upper tail, its correlations were interpreted alongside group profiles and influence diagnostics \cite{deWinter2016}. The analysis also checked support for joint household categories and compared unweighted and expansion-weighted means. Display caps pooled sparse upper categories in figures and cross-classification; regression retained the original eligible counts.

The eligible sample covered 1,179 traffic zones, including 1,158 represented in training and 1,081 in testing. Because households from a zone could occur in both partitions, the test measures prediction within the 2006 survey setting. The fixed test partition remained outside coefficient estimation, information-criterion ranking, and estimation of cross-classification rates. Earlier workflow development involved exploration of this dataset; the resulting test assessment is therefore an internal validation exercise.

\subsection{Multiple linear regression}
\begin{samepage}
Multiple linear regression was used to describe the conditional association between household characteristics and the expected number of recorded trips. For household $i$, the working model was
\begin{equation}
 y_i = \beta_0 + \sum_{j=1}^{p-1} x_{ij}\beta_j + \varepsilon_i,
 \qquad \widehat y_i = \widehat\beta_0 + \sum_{j=1}^{p-1}x_{ij}\widehat\beta_j,
 \label{eq:mlr}
\end{equation}
\end{samepage}
where $x_{ij}$ is a predictor or constructed predictor term, $p$ includes the intercept, and $\varepsilon_i$ represents variation not captured by the specification. The model is linear in its coefficients when predictors are multiplied to represent interactions. Its predictions are conditional mean trip rates and need not be integers.

Let $P_i$, $V_i$, $L_i$, $F_i$, $Q_i$, $H_i$, and $S_i$ denote household persons, vehicles, licensed drivers, full-time workers, part-time workers, home workers, and students, respectively. Household size, vehicle count, and driver count were centered as $P_{ci}=P_i-3$, $V_{ci}=V_i-1$, and $L_{ci}=L_i-2$. In specifications with interactions, the coefficient of a centered count is a local slope at the reference values of the interacting counts. A one-unit change is evaluated from the full equation. These fixed reference values were chosen for this project. Centering leaves the fitted values of the corresponding complete specification unchanged.

The starting specification was
\begin{equation}
 y_i=\beta_0+\beta_P P_{ci}+\beta_V V_{ci}
       +\beta_F F_i+\beta_S S_i+\varepsilon_i.
 \label{eq:simple-model}
\end{equation}
Persons, workers, and students represent potential participation in daily activities, while vehicles represent access to motorized travel. The coefficients describe conditional associations within the observed households.

The candidate set extended the starting model to licensed drivers, alternative employment measures, dwelling type, region, and two interactions. Driver count indicates capacity for independent motorized travel; home working may change commuting and other travel. Dwelling type and region summarize residential context. Interactions let the driver and vehicle associations vary with household size.

The searched effects were selected from the following encompassing specification:
\begin{align}
 y_i={}&\beta_0+\beta_P P_{ci}+\beta_V V_{ci}+\beta_L L_{ci}
       +\beta_F F_i+\beta_Q Q_i+\beta_H H_i+\beta_S S_i \notag\\
      &+\boldsymbol{\gamma}^{\mathsf T}\mathbf d_i
       +\boldsymbol{\eta}^{\mathsf T}\mathbf r_i \notag\\
      &+\beta_{PL}P_{ci}L_{ci}+\beta_{PV}P_{ci}V_{ci}
       +\varepsilon_i.
 \label{eq:alternative-model}
\end{align}
Here $\mathbf d_i$ and $\mathbf r_i$ contain indicator variables for dwelling type and region. Their reference categories were House and Toronto. The other regions were Durham, York, and Peel, assigned from residential zone ranges. The search treated total workers, $F_i+Q_i$, as a substitute for the separate employment counts, avoiding redundant effects. It used the region categories to summarize location, while household identifiers, exact zone identifiers, survey administration fields, and the expansion factor remained outside the predictor set.

\subsubsection{Estimation and specification search}
All coefficients were estimated by unweighted ordinary least squares on the training households:
\begin{equation}
 \widehat{\boldsymbol\beta}
 =\underset{\boldsymbol\beta}{\operatorname{arg\,min}}
   \sum_{i\in\mathcal T}(y_i-\mathbf x_i^{\mathsf T}\boldsymbol\beta)^2,
 \label{eq:ols}
\end{equation}
where $\mathcal T$ denotes the training set and $\mathbf x_i$ includes the intercept. This gives each sampled household equal weight in fitting. The same training observations were used for every candidate model.

The search evaluated 1,040 valid effect subsets, including the intercept-only specification; household size was optional. Categorical variables entered as whole factors, and interactions required their supporting main effects. Candidates were scored using working Gaussian OLS likelihoods:
\begin{equation}
 \mathrm{AIC}=-2\widehat\ell+2p,
 \qquad
 \mathrm{BIC}=-2\widehat\ell+p\log(n_{\mathrm{tr}}),
 \label{eq:ic}
\end{equation}
where $n_{\mathrm{tr}}=|\mathcal T|$ and $\widehat\ell$ is the maximized working log-likelihood. The penalty counts the regression coefficients, including the intercept, in $p$; a common variance-parameter penalty leaves the ranking unchanged. The working Gaussian likelihood provided a comparison rule, while the observed trip-count distribution informed the residual and prediction checks.

Models were examined in ascending BIC order. The selected alternative was the first model for which every removable whole effect had a zone-clustered Wald test $p$-value below 0.01. Only main effects required by retained interactions were protected from removal. The AIC leaders and the best additive specification were retained for comparison. The squared terms were removed after the first evaluation because their parabolic form made predictions fall at high vehicle counts and produced negative predictions. Survey weekday was subsequently removed from the candidate effects to simplify the specification. The revised search used training data; its test results constitute internal validation after this specification revision.

Coefficient uncertainty was estimated using a sandwich covariance clustered by residential traffic zone, with the finite-sample correction
$G/(G-1)\,(n_{\mathrm{tr}}-1)/(n_{\mathrm{tr}}-p)$, where $G=1,158$ training zones. Coefficient intervals used a $t$ reference distribution with $G-1$ degrees of freedom; whole-factor tests used the corresponding clustered Wald $F$ statistic. This permits correlated errors within zones and unequal error variances, while retaining the OLS coefficient estimates. The reported 99\% intervals and tests for the selected alternative are exploratory because they do not adjust for searching among specifications.

Training goodness of fit was summarized by
\begin{equation}
 R^2=1-\frac{\sum_{i\in\mathcal T}(y_i-\widehat y_i)^2}
                 {\sum_{i\in\mathcal T}(y_i-\overline y_{\mathcal T})^2},
 \qquad
 \overline R^2=1-(1-R^2)\frac{n_{\mathrm{tr}}-1}{n_{\mathrm{tr}}-p}.
\end{equation}
Residual plots, a normal quantile plot, residual skewness and kurtosis, variance-inflation factors, leverage, Cook's distance, and externally studentized residuals were examined. Leverage above $3p/n_{\mathrm{tr}}$, Cook's distance above $4/n_{\mathrm{tr}}$, and absolute studentized residuals above three flagged observations for review. The robust Breusch--Pagan diagnostic provided an independent-observation screen for heteroscedasticity. Predicted trip curves at fixed household profiles were examined across the observed vehicle range together with exact-count training support. Negative predictions remained in the evaluation so that errors represented the fitted OLS model.

\subsection{Cross-classification}
The cross-classification model assigned an average trip rate to households sharing a size and licensed-driver category, an established trip-generation approach \cite{Moeckel2017}. The primary size groups were 1, 2, 3, 4, and 5 or more persons; driver groups were 0, 1, 2, and 3 or more. Household size represents potential activity demand, while driver availability describes one aspect of mobility access. Their joint categories allow mean trip rates to vary across combinations. Pooling the upper categories increases sample support in the tails.

For size group $a$ and driver group $b$, let $\mathcal T_{ab}$ be the corresponding training households and $n_{ab}=|\mathcal T_{ab}|$. The cell rate was the unweighted mean
\begin{equation}
 \widehat r_{ab}=\frac{1}{n_{ab}}\sum_{i\in\mathcal T_{ab}}y_i.
\end{equation}
A feasible cell with fewer than 30 training households used the training mean for its size group, $\widehat r_a$. Thus, the prediction for a feasible test household was
\begin{equation}
 \widehat y_i^{\mathrm{CC}}=
 \begin{cases}
 \widehat r_{ab}, & n_{ab}\geq30,\\
 \widehat r_a, & n_{ab}<30.
 \end{cases}
\end{equation}
The threshold of 30 set a minimum cell count for direct rate estimation. Precision still varied with cell size and within-cell variation. Structurally impossible combinations, such as two licensed drivers in a one-person household, had blank rates and heatmap cells. Sparse-cell counts and fallback predictions applied to feasible combinations. Tables and the heatmap reported the household count supporting each observed rate.

Two alternative boundary schemes were evaluated: size capped at 4 or more with drivers capped at 3 or more, and size capped at 5 or more with drivers capped at 4 or more. All rates and fallback values were fitted on training households. Test results for these alternatives served as boundary sensitivity checks. The primary boundaries were fixed from the training-data support assessment.

\subsection{Household prediction assessment}
Regression and cross-classification were evaluated on the same $n_{\mathrm{te}}=27,978$ eligible test households. A benchmark assigned every test household the unweighted training mean, $\widehat y_i^{\mathrm{mean}}=\overline y_{\mathcal T}$. A regression using only household size and licensed drivers gave a common-predictor comparison with cross-classification; it retained exact counts where cross-classification grouped upper counts. The simple four-predictor regression, best additive alternative, and selected regression showed how performance changed across broader specifications. Both principal methods were also assessed within the same size and driver groups.

For prediction error $e_i=\widehat y_i-y_i$, the household measures were
\begin{align}
 \mathrm{MAE}&=\frac{1}{n_{\mathrm{te}}}\sum_{i\in\mathcal E}|e_i|,\\
 \mathrm{RMSE}&=\sqrt{\frac{1}{n_{\mathrm{te}}}\sum_{i\in\mathcal E}e_i^2},\\
 \mathrm{ME}&=\frac{1}{n_{\mathrm{te}}}\sum_{i\in\mathcal E}e_i,
\end{align}
where $\mathcal E$ is the test set. MAE gives the average absolute error in trips per household. RMSE gives greater weight to large errors. ME is the signed mean error in trips: a positive value indicates overprediction. A small ME can coexist with substantial household errors because positive and negative errors cancel. Observed-versus-predicted plots and results for selected household groups supplemented these overall measures.

The expansion factor $w_i=\texttt{expf}_i$ was used for an additional weighted assessment, without re-estimating either model. Weighted MAE and ME replaced $1/n_{\mathrm{te}}$ by normalized weights $w_i/\sum_{j\in\mathcal E}w_j$; weighted RMSE was
\begin{equation}
 \mathrm{RMSE}_{w}=\sqrt{\frac{\sum_{i\in\mathcal E}w_i e_i^2}
                                  {\sum_{i\in\mathcal E}w_i}}.
\end{equation}

\subsection{Traffic-zone aggregation and scope of inference}
For each zone $z$ represented in the test set, observed and predicted sample totals were summed over the same records $\mathcal E_z$:
\begin{equation}
 Y_z=\sum_{i\in\mathcal E_z}y_i,
 \qquad \widehat Y_z=\sum_{i\in\mathcal E_z}\widehat y_i.
\end{equation}
Expansion-weighted sums were calculated consistently as
\begin{equation}
 Y_z^{w}=\sum_{i\in\mathcal E_z}w_i y_i,
 \qquad \widehat Y_z^{w}=\sum_{i\in\mathcal E_z}w_i\widehat y_i.
\end{equation}
Zone MAE, RMSE, and ME used the corresponding differences, with each of the 1,081 evaluated zones receiving equal weight in these summaries. Zones without test households were not included. These measures and zone agreement plots assess how household errors accumulate or cancel when aggregated.

The weighted sums represent the contributions of the test records under their original survey expansion factors. They are sample-based zone aggregates because the analysis used neither an inverse test-selection adjustment nor population recalibration. Both model comparisons use the same records and weighting convention. The results describe associations and predictive performance within this historical survey sample. Forecasts for another period or geography would require separate validation.
